\section{Proper pro-three boundary fibres of Fermat curves}
\label{sec:finitefield3-boundary-reconstruction}

\subsection{Exact scope and results}
\label{subsec:ff3-scope}

We determine the possible identifications among rational coordinate
boundary points in the proper geometric pro-three group.
Let \(n\ge3\), let \(k=\mathbb F_Q\) have characteristic prime to \(3n\),
and let
\[
 C_n:\ X^n+Y^n+Z^n=0.
\]
This smooth proper curve has genus \((n-1)(n-2)/2\ge1\). Put
\[
 D_X=C_n\cap\{X=0\},\quad D_Y=C_n\cap\{Y=0\},\quad
 D_Z=C_n\cap\{Z=0\},\qquad D=D_X\sqcup D_Y\sqcup D_Z.
\]
Fix a geometric base point and write
\[
 N=\pi_1(C_{n,\bar k})^{\mathrm{pro}\text{-}3},\qquad
 \Pi=\pi_1(C_n)/\ker\bigl(\pi_1(C_{n,\bar k})\longrightarrow N\bigr).
\]
Thus
\begin{equation}
 1\longrightarrow N\longrightarrow\Pi\longrightarrow G_k
 \longrightarrow1,\qquad G_k\simeq\widehat{\mathbb Z}.
 \label{eq:ff3-1-1}
\end{equation}
The arithmetic quotient is full \(G_k\). Let \(\kappa(P)\) be the
proper point-section class of \(P\in C_n(k)\), modulo conjugacy by \(N\).
All fibres below are intersected with \(D(k)\); no assertion excludes
an ordinary proper point from the full fibre of a boundary point.

\begin{theorem}[Separation of coordinate components]
\label{thm:ff3-different-components}
If \(P,Q\in D(k)\) lie on different coordinate components, then
\begin{equation}
 \kappa(P)\ne\kappa(Q).
 \label{eq:ff3-1-2}
\end{equation}
This holds for both odd and even \(n\).
\end{theorem}

\begin{theorem}[Uniform boundary blocks and coefficient-prime separation]
\label{thm:ff3-dichotomy}
Suppose \(D(k)\ne\varnothing\), and put \(d=\gcd(n,Q-1)\).
Each component has \(d\) rational points. There is one divisor \(h\mid d\)
such that the fibres on each component are the cosets of the unique
subgroup \(\mu_h(k)\subset\mu_d(k)\) under its regular diagonal action.
There are exactly \(3d/h\) boundary section classes.
Moreover,
\[
 3\nmid h,\qquad 3\mid n\ \Longrightarrow\ h=1.
\]
In particular, the boundary map is injective for every exponent
divisible by three, at every allowed finite field.

For a prime exponent \(a\ge5\), if \(Q\not\equiv1\pmod a\), the
boundary consists of three points with distinct classes. If
\(Q\equiv1\pmod a\), either all \(3a\) classes are distinct or the
three whole coordinate components are exactly the fibres.
\end{theorem}

For odd \(n\), the point \((0:-1:1)\) shows that \(D(k)\ne\varnothing\).
For even \(n\), nonemptiness is equivalent to
\(-1\in(k^\times)^n\), or to \((Q-1)/d\) being even.
When this condition fails, all three rational boundary components
are empty.

\begin{corollary}[Specified primes with singleton boundary fibres]
\label{cor:ff3-density}
Let \(n\ge3\) be odd. If \(3\mid n\), the boundary map is injective
at every rational prime \(q\nmid3n\). If \(3\nmid n\), the primes
\(q\nmid3n\) satisfying \(\gcd(n,q-1)=1\) give singleton boundary
fibres. This specified sufficient set has natural density
\begin{equation}
 \prod_{p\mid n}\frac{p-2}{p-1}.
 \label{eq:ff3-1-3}
\end{equation}
In particular the specified nonsplit sets for \(n=5\) and \(n=7\)
have densities \(3/4\) and \(5/6\), respectively.
\end{corollary}

The proof first uses Pál's order-three action theorem to separate
different components. A tangent-line calculation then proves the
stronger injectivity when three divides the exponent. The remaining
same-component alternative for prime exponents different from three
is not decided; in particular no density-one separation or affine
global classification at the base exponents five and seven follows.

\subsection{Geometric centralizers over a finite field}
\label{subsec:ff3-centralizers}

We need uniqueness, not just existence, of invariant geometric paths. The exact published input is \citet[Lemma~1.7, p.~144]{FCSaidiTamagawa2009}. Its standing assumptions on p.~140 allow any retained prime set containing a prime different from the residue characteristic, and allow a proper hyperbolic curve. Thus \(\Sigma=\{3\}\) is allowed when the genus is at least two. The direct proper-curve proof below also covers genus one. The additional cofinite-\(\Sigma\) assumptions occurring later in their paper are not present in this lemma.

\begin{lemma}[Trivial geometric centralizers]
\label{lem:ff3-centralizer}

Let \(C/k\) be a smooth proper geometrically connected curve of positive genus over a finite field of characteristic different from \(\ell\). Put \[
 N_C=\pi_1(C_{\bar k})^{\mathrm{pro}\text{-}\ell},
\] and retain full \(G_k\) in the corresponding arithmetic quotient. For any arithmetic section \(s\), \begin{equation}
Z_{N_C}(s(G_k))=1.
\label{eq:ff3-2-1}
\end{equation}

\end{lemma}
\begin{proof}

In genus at least two, apply Saïdi--Tamagawa Lemma 1.7 to \(Z=s(G_k)\). Projection is injective on \(Z\) and its image is all of \(G_k\), hence open. The lemma injects the commensurator of \(Z\), and therefore its centralizer, into \(G_k\). Its intersection with the geometric kernel is consequently trivial.

For every positive genus, the following proper, singleton-prime argument proves the assertion directly. Use \(s\) to identify the arithmetic quotient with \(N_C\rtimes G_k\), and suppose a nonidentity \(z\in N_C\) is fixed by the \(G_k\)-action. Choose an open characteristic subgroup \(M\subset N_C\) with \(z\notin M\), and put \[
 H=\overline{\langle M,z\rangle}.
\] This is an open \(G_k\)-stable subgroup. The image of \(z\) in \(H/M\) is nonzero, and \(H/M\) is cyclic; hence the image of \(z\) in \(H^{\mathrm{ab}}\) is nonzero.

The open subgroup \(H\rtimes G_k\) corresponds to a finite étale geometrically connected proper cover \(Y\to C\) over \(k\). The geometric pro-\(\ell\) topology induced on its geometric subgroup is the full topology: an open subgroup of a pro-\(\ell\) group has this property. At the original fundamental-group level, one proves it by first passing to the normal core of the covering subgroup, then taking the intersection of the finitely many conjugates of the kernel of a finite \(\ell\)-quotient. The resulting quotient is an extension of finite \(\ell\)-groups. Thus \[
 H=\pi_1(Y_{\bar k})^{\mathrm{pro}\text{-}\ell},
 \qquad H^{\mathrm{ab}}=T_\ell\operatorname{Jac}(Y).
\] Since \(Y\) is proper, \(\operatorname{Jac}(Y)(k)\) is finite. Its rational \(\ell\)-primary torsion group is finite and therefore admits no nonzero compatible Tate-module system, so \[
 (T_\ell\operatorname{Jac}(Y))^{G_k}=0.
\] But the nonzero image of \(z\) is fixed by \(G_k\), a contradiction. This proves \eqref{eq:ff3-2-1}.

\end{proof}

This proper proof avoids any need to infer integral fixed-point vanishing from a rational vector-space statement without checking torsion. The cited original lemma proves the more general assertion by the same open-subgroup device and Frobenius weights.

\subsection{Arithmetic equality as an invariant-path statement}
\label{subsec:ff3-paths}

Take \(b\in C(k)\), and use its point section to fix the actual \(G_k\)-action on \[
 N_C=\pi_1(C_{\bar k},b)^{\mathrm{pro}\text{-}\ell}.
\] For geometric points \(x,y\), let \(\mathcal P_\ell(x,y)\) be the torsor of pro-\(\ell\) étale paths between their fibre functors.

If \(x,y\) are rational, this torsor has its natural \(G_k\)-action. Choosing a path identifies its descent cocycle with the difference between the two point sections. Therefore \begin{equation}
\kappa(x)=\kappa(y)
 \quad\Longleftrightarrow\quad
 \mathcal P_\ell(x,y)^{G_k}\ne\varnothing.
\label{eq:ff3-3-1}
\end{equation} This is the usual torsor description of nonabelian \(H^1\), with no change of the arithmetic quotient. Two invariant paths differ by a \(G_k\)-fixed element of the appropriate geometric fundamental group. Lemma~\ref{lem:ff3-centralizer} makes that element the identity. Thus, in our setting, an invariant path is unique whenever it exists.

Now let \(A\) be a finite group of \(k\)-automorphisms of \(C\). Define a geometric action extension \begin{equation}
1\longrightarrow N_C\longrightarrow E_A
 \overset{\mathrm{pr}}\longrightarrow A\longrightarrow1.
\label{eq:ff3-3-2}
\end{equation} One concrete construction uses pairs \((a,p)\), with \(a\in A\) and \(p\in\mathcal P_\ell(b,a(b))\). Multiplication concatenates \(p\) with the image under \(a\) of the second path. This construction is independent, up to the usual inner equivalence, of the base-point choices.

The quotient \(A\) in \eqref{eq:ff3-3-2} is retained in its entirety. Only the geometric kernel has been made pro-\(\ell\). In particular, for \(A=S_3\) and \(\ell=3\), \(E_A\) is \textbf{not} being defined as the maximal pro-three quotient of an orbifold group; doing that would lose its quotient of order two.

Because the automorphisms are defined over \(k\) and \(b\) is rational, \(G_k\) acts on \(E_A\), preserves its group law, and acts trivially on its quotient \(A\).

\begin{lemma}[Coherent lifting of an invariant point class]
\label{lem:ff3-invariant-action}

Suppose \(C/k\) satisfies Lemma~\ref{lem:ff3-centralizer}. If \(\kappa(b)\) is fixed by the action of every element of \(A\), then \eqref{eq:ff3-3-2} has a section \[
 A\longrightarrow E_A
\] whose image is pointwise fixed by \(G_k\).

\end{lemma}
\begin{proof}

For each \(a\in A\), invariance says \(\kappa(a(b))=\kappa(b)\). By \eqref{eq:ff3-3-1}, the fibre of \(E_A\to A\) above \(a\) contains a \(G_k\)-fixed element. Its uniqueness follows from Lemma~\ref{lem:ff3-centralizer}. Therefore \[
 E_A^{G_k}\longrightarrow A
\] is bijective. It is a homomorphism because \(G_k\) acts on \(E_A\) by automorphisms. Its inverse is the required section.

\end{proof}

The uniqueness step is essential. Choosing unrelated paths for each automorphism would not by itself prove a finite-group splitting.

\subsection{Pál’s theorem and its admissible application}
\label{subsec:ff3-pal}

Let \(C/\mathbb C\) be a smooth projective connected curve of positive genus, and let a cyclic group \(B\) of prime order \(\ell\) act nontrivially. \citet[Theorem~5.2, author pp.~12--13]{FCPal2014} states that the fixed points of \(B\) correspond bijectively to conjugacy classes of sections of the pro-\(\Sigma\) completed action extension, when \(\ell\in\Sigma\). For \(\Sigma=\{\ell\}\) this is the full pro-\(\ell\) action extension.

Pál's definition retains the finite quotient only when every prime dividing its order belongs to \(\Sigma\). Our only use of this theorem is \begin{equation}
B=C_3,\qquad \Sigma=\{3\}.
\label{eq:ff3-4-1}
\end{equation} We do not apply it to \(S_3\) or to a group of order \(5\) or \(7\).

For \(B=C_3\), the path extension \(E_B\) of Section~\ref{subsec:ff3-paths} is the same extension as in Pál's theorem. Its kernel is pro-three and its quotient has order three, so the total group is pro-three. Equivalently, Pál's exact completion sequence (Definition 2.3) identifies the pro-three completion of the topological action extension with this geometric path extension.

There is an elementary naturality that we will need. If \(\tau\) normalizes \(B\), it acts on section classes of \(E_B\) by conjugating the finite quotient and transporting paths. The fixed-point-to-section bijection is equivariant: \begin{equation}
\tau\cdot[s_x]=[s_{\tau x}].
\label{eq:ff3-4-2}
\end{equation} Indeed, in Pál's Definition 2.1, the subgroup associated with \(x\) is the subgroup of lifts fixing a chosen point above \(x\) in the universal cover. Applying a normalizing automorphism sends it to the subgroup fixing the transported point over \(\tau x\). Changing the chosen lift only conjugates by the geometric group.

\paragraph{Specialization to finite characteristic}

The Fermat model and its coordinate permutations are defined over \(\mathbb Z[1/n]\). Adjoin a primitive cube root \(\omega\), and use a strict henselian discrete valuation ring above the residue characteristic, with \(3n\) invertible. The Fermat curve is smooth and proper over this ring. When \(3\nmid n\), the two points \begin{equation}
f_+=(1:\omega:\omega^2),\qquad
 f_-=(1:\omega^2:\omega)
\label{eq:ff3-4-3}
\end{equation} are sections over it.

Smooth proper specialization identifies the geometric pro-three fundamental groups and their path torsors, compatibly with morphisms and composition. The exact original references are \citet[Exposé~X, Section~2, pp.~206--207; Theorem~3.8 and Corollary~3.9, p.~217 (original marginal pp.~283--284)]{SGA1}. The generic geometric fibre can be compared with the complex curve by characteristic-zero base change and Riemann existence \citep[Exposé~XII, Theorem~5.1]{SGA1}.

The needed naturality follows from the specialization construction on finite étale covers and their fibre functors: pullback by each extending coordinate automorphism commutes with restriction to the geometric generic and special fibres, and evaluation at an extending point commutes with those restriction functors. Consequently the fixed-point splitting classes and the normalizer action in \eqref{eq:ff3-4-2} specialize together.

Only a \textbf{geometric} comparison is asserted. The arithmetic Galois groups in characteristic zero and finite characteristic are not identified. Nor is a specialization theorem for a possibly wild \(S_3\)-quotient stack required.

This also covers residue characteristic \(2\). The geometric coefficient prime and the order of \(C_3\) are still invertible; the transposition is used only as an extending automorphism, not as a tame group quotient of order two.

\subsection{The coordinate permutation extension does not split}
\label{subsec:ff3-s3}

We apply the finite-action theorem to the cyclic subgroup of coordinate
permutations. On \(C_n\), put
\[
 \rho(X:Y:Z)=(Y:Z:X),\qquad \tau(X:Y:Z)=(X:Z:Y).
\]
They act faithfully and generate \(S_3\), with
\(\rho^3=\tau^2=1\) and \(\tau\rho\tau=\rho^{-1}\).
A projective fixed point of \(\rho\) must be one of its three eigenlines,
\[
 (1:1:1),\quad (1:\omega:\omega^2),\quad(1:\omega^2:\omega).
\]
The first has Fermat value \(3\ne0\). If \(3\mid n\), the other two
also have Fermat value three, so the cyclic fixed locus is empty.
If \(3\nmid n\), precisely those two lie on the curve, and \(\tau\)
interchanges them.

\begin{lemma}[The coordinate action extension does not split]
\label{lem:ff3-s3-nonsplit}
The relative geometric action extension
\begin{equation}
 1\longrightarrow N\longrightarrow E_{S_3}
 \longrightarrow S_3\longrightarrow1
 \label{eq:ff3-5-1}
\end{equation}
has no splitting.
\end{lemma}
\begin{proof}
Suppose \(t:S_3\to E_{S_3}\) were a splitting, and restrict it to
\(B=\langle\rho\rangle\), obtaining \(u:B\to E_B\).
If \(3\mid n\), Pál's theorem and the specialization just established
identify its section classes with an empty fixed locus, a contradiction.

If \(3\nmid n\), there are exactly two possible classes, those of
the two cyclic fixed points. The transposition interchanges them.
On the other hand,
\[
 t(\tau)\,u(\tau^{-1}g\tau)\,t(\tau)^{-1}=u(g)\qquad(g\in B).
\]
Changing the chosen lift of \(\tau\) changes this action only by
geometric conjugacy. Thus the class of \(u\) would be fixed by the
transposition, again a contradiction.
\end{proof}

Only the cyclic order-three subgroup is an input to Pál's theorem.
The finite quotient \(S_3\) in \eqref{eq:ff3-5-1} has not been
replaced by a pro-three completion.

\subsection{Separation over both odd and even exponents}
\label{subsec:ff3-canonical}

Twisted transpositions compare a pair of components without assuming
that the exponent is odd. After permuting coordinates, write
\[
 P=(0:u:1),\qquad Q=(v:0:1),\qquad u^n=v^n=-1,
\]
with \(u,v\in k^\times\). Let
\[
 A=\operatorname{diag}(1,u^2,1),\qquad
 B=\operatorname{diag}(v^2,1,1),\qquad \sigma=(XZ).
\]
These diagonal maps preserve \(C_n\), since \(u^{2n}=v^{2n}=1\).
They satisfy
\[
 A(Q)=Q,\quad A\tau(P)=(0:u^2:u)=P,\qquad
 B(P)=P,\quad B\sigma(Q)=(v^2:0:v)=Q
\]
as projective points.
If \(\kappa(P)=\kappa(Q)=c\), both \(A\) and \(A\tau\) preserve \(c\),
so \(\tau\) preserves \(c\). Likewise \(\sigma\) preserves \(c\).
The two transpositions generate \(S_3\).
Lemma~\ref{lem:ff3-invariant-action} would split
\eqref{eq:ff3-5-1}, contradicting
Lemma~\ref{lem:ff3-s3-nonsplit}. This proves
Theorem~\ref{thm:ff3-different-components}. For odd \(n\), in particular,
\begin{equation}
 \kappa(0:-1:1),\quad\kappa(-1:0:1),\quad\kappa(1:-1:0)
 \quad\hbox{are pairwise distinct}.
 \label{eq:ff3-6-1}
\end{equation}
The automorphisms used are defined over \(k\). The cyclic fixed points
used to obstruct a geometric splitting may be considered over \(\bar k\);
no cube root of unity in \(k\) is required.

\subsection{Uniform cyclic blocks}
\label{subsec:ff3-components}

The diagonal symmetries determine the form of every possible fibre.
If \(u^n=-1\) has one solution in \(k\), all its solutions are
\(u\zeta\) for \(\zeta\in\mu_n(k)=\mu_d(k)\), where \(d=\gcd(n,Q-1)\).
Thus each coordinate component is a regular orbit of this cyclic group.
The relation
\[
 P\sim Q\quad\Longleftrightarrow\quad\kappa(P)=\kappa(Q)
\]
is invariant under every \(k\)-automorphism. The elements of the cyclic
group carrying a fixed point to an equivalent point form a subgroup:
if \(tP\sim P\) and \(uP\sim P\), invariance gives \(tu^{-1}P\sim P\).
The fibres are its cosets. A cyclic group has a unique subgroup of
each order, so their common size is a divisor \(h\mid d\).

Coordinate permutations identify these relations on the three
components, possibly inverting the cyclic coordinate. They preserve
the subgroup order, so the same \(h\) works on all three.
Theorem~\ref{thm:ff3-different-components} excludes any identification
between components. There are exactly \(3d/h\) classes.

If \(n=a\ge5\) is prime, then \(d=1\) or \(a\).
In the first case the three boundary points are distinct; in the
second, primality leaves \(h=1\) or \(a\), giving the asserted
singleton or whole-component dichotomy. We prove the stronger
conclusion \(h=1\) when \(3\mid n\) next.

\subsection{Boundary divisor classes and the coefficient prime}
\label{subsec:ff3-abelian}

An explicit pole calculation distinguishes the exponents divisible
by three from the remaining ones.

\begin{lemma}[Exact order of a coordinate-point difference]
\label{lem:ff3-exact-cusp-order}
Over an algebraically closed field of characteristic not dividing
\(n\ge3\), the divisor class \(P-Q\) of any two distinct coordinate
points on \(C_n\) has exact order \(n\) in its Jacobian.
\end{lemma}
\begin{proof}
At a coordinate point the tangent line intersects the curve in
\(n\) times that point. Therefore the ratio of the tangent lines at
\(P,Q\) has divisor \(n(P-Q)\).

It remains to exclude a proper divisor of \(n\) as the order.
After permuting coordinates, write \(P=(0:u:1)\), with \(u^n=-1\).
The tangent is \(L=Y-uZ\), and its zero divisor on the curve is \(nP\).
On \(C_n\setminus\{P\}\), use affine coordinates
\[
 x=X/L,\qquad z=Z/L.
\]
Their coordinate ring is
\[
 \bar k[x,z]\big/\bigl(x^n+(1+uz)^n+z^n\bigr).
\]
The polynomial in \(z\) has degree \(n-1\), since \(u^n=-1\)
and its leading coefficient is \(n u^{n-1}\ne0\).
At \(P\), \(X/Z\) is a uniformizer, \(L\) vanishes to order \(n\),
and \(Z\) is a unit. Hence \(x\) has pole order \(n-1\), and \(z\)
has pole order \(n\).

Every function regular away from \(P\) has a unique expression
\[
 \sum_{i=0}^{n-1} a_i(z)x^i,\qquad a_i(z)\in\bar k[z],
\]
because the equation is monic of degree \(n\) in \(x\).
The monomials \(z^j x^i\), \(j\ge0\), \(0\le i<n\), have pole orders
\(jn+i(n-1)\). These orders are pairwise distinct: reduction modulo
\(n\) determines \(i\), and then determines \(j\).
Thus leading poles cannot cancel, and every nonconstant function
whose only possible pole is \(P\) has pole order at least \(n-1\).

If the order of \(P-Q\) were a proper divisor \(m\) of \(n\),
a function with divisor \(m(Q-P)\) would have its only pole at
\(P\), of order \(m\le n/2<n-1\). This is impossible.
\end{proof}

Suppose now \(P,Q\in D(k)\) are distinct and \(3\mid n\).
If their proper pro-three point sections were equal, their images
in the geometric abelianization \(T_3\operatorname{Jac}(C_n)\)
would have the same Kummer class. The finite Kummer exact sequences
would then give
\[
 [P-Q]\in\bigcap_{j\ge1}3^j\operatorname{Jac}(C_n)(k).
\]
The group \(\operatorname{Jac}(C_n)(k)\) is finite; that intersection
is its prime-to-three subgroup. This contradicts the exact order
\(n\) in Lemma~\ref{lem:ff3-exact-cusp-order}.
Consequently every boundary fibre is a singleton when \(3\mid n\).
If \(3\nmid n\), then \(3\nmid d\), so the block size is automatically
prime to three. This proves all of Theorem~\ref{thm:ff3-dichotomy}.

When \(3\nmid n\), the same divisor calculation shows that every
boundary difference has zero three-primary Kummer image. The
original Jacobian therefore does not separate these boundary points.
The component-separation theorem in that range uses the full
nonabelian proper group. Distinguishing affine cusp inertia alone
would not prove separation after that inertia is filled.

For the density assertion, if \(3\mid n\), injectivity holds at every
good prime just proved. Otherwise, for odd \(n\), the condition
\(\gcd(n,q-1)=1\) gives precisely the three rational coordinate
points. For each distinct prime \(p\mid n\), it excludes the identity
class among the \(p-1\) invertible residues modulo \(p\).
Chinese remaindering and the cyclotomic form of Chebotarev give
\eqref{eq:ff3-1-3}; this is natural density, hence also Dirichlet density
\citep[Definition~8.30, Theorem~8.31, and Example~8.34,
pp.~147--148]{Milne2020ANT}.
This proves Corollary~\ref{cor:ff3-density} and
Theorem~\ref{thm:finitefield-boundary-intro}.

Each proved inequality of full point-section classes is visible in
some finite characteristic pro-three geometric quotient. Otherwise
the closed sets of conjugators giving equality in finite quotients
would have the finite-intersection property, and compactness would
give equality in the full group. No uniform small quotient is asserted.

\subsection{A sufficient automorphism criterion for collisions}
\label{subsec:ff3-automorphisms}

The following criterion is useful for testing potential constructions. It is a sufficient condition only.

\begin{lemma}[The Frattini automorphism kernel]
\label{lem:ff3-frattini}

For a finitely generated pro-\(\ell\) group \(M\), \[
 \ker\bigl(\operatorname{Aut}(M)\to
            \operatorname{Aut}(M/\Phi(M))\bigr)
\] is pro-\(\ell\), where \(\Phi(M)=\overline{M^\ell[M,M]}\). The corresponding kernel in the outer automorphism group is also pro-\(\ell\).

\end{lemma}
\begin{proof}

For a finite \(\ell\)-group \(P\), choose an ordered basis of \(P/\Phi(P)\), of dimension \(d\). It has \(|\Phi(P)|^d\) ordered lifts to \(P\). Every lift generates \(P\), by Burnside's basis theorem. An automorphism trivial on \(P/\Phi(P)\) acts on the set of lifts, and the stabilizer of any lift is trivial: an automorphism fixing a generating tuple is the identity. The order of the automorphism kernel therefore divides the \(\ell\)-power \(|\Phi(P)|^d\), so it is an \(\ell\)-group.

Pass to a cofinal family of characteristic finite quotients of \(M\). The inverse-limit kernel is pro-\(\ell\). Inner automorphisms act trivially on the Frattini quotient, so the outer kernel is a quotient of this pro-\(\ell\) kernel.

\end{proof}

This is a statement about the kernel. The full automorphism group is not asserted to be pro-\(\ell\).

\begin{lemma}[A residual-trivial automorphism criterion]
\label{lem:ff3-trivial-residual}

Let \(C/k\) be a geometrically connected proper curve with \(b\in C(k)\) and finitely generated geometric pro-\(\ell\) group \(M\). Let \(\alpha\in\operatorname{Aut}_k(C)\) have finite order \(r\) prime to \(\ell\). If \(\alpha\) acts trivially on \(M/\Phi(M)\), then \begin{equation}
\kappa(\alpha b)=\kappa(b).
\label{eq:ff3-9-1}
\end{equation}

\end{lemma}
\begin{proof}

The outer action of \(\alpha\) has order dividing \(r\) and lies in the pro-\(\ell\) kernel of Lemma~\ref{lem:ff3-frattini}. It is therefore trivial.

Form the geometric action extension \(1\to M\to E_{\langle\alpha\rangle}\to C_r\to1\). Since every element of \(C_r\) acts on \(M\) by an inner automorphism, the centralizer \[
 L=Z_{E_{\langle\alpha\rangle}}(M)
\] surjects onto \(C_r\), with kernel \(Z(M)\). This is a central extension of \(C_r\) by the abelian pro-\(\ell\) group \(Z(M)\).

Choose a lift \(e\in L\) of a generator of \(C_r\). Then \(e^r\in Z(M)\). Multiplication by \(r\) is an automorphism of an abelian pro-\(\ell\) group. There is a unique \(z\in Z(M)\) with \(z^r=e^{-r}\), so \(t=ze\) has order exactly \(r\) and maps to that generator. The subgroup \(\langle t\rangle\) is the unique complement to \(Z(M)\) in \(L\): any other lift of the same generator differs by an element of \(Z(M)\), and its \(r\)-th power forces that element to be the identity.

The \(G_k\)-action preserves \(L\), its kernel, and this unique complement. It acts trivially on the quotient \(C_r\) because \(\alpha\) is defined over \(k\). Hence it fixes the complement pointwise. In particular it fixes a path over \(\alpha\), and \eqref{eq:ff3-3-1} gives \eqref{eq:ff3-9-1}.

\end{proof}

No centre-freeness of \(M\) is required in this proof. The central kernel is retained and the complement is constructed explicitly.

\paragraph{Failure of this criterion for the Fermat diagonal rotations}

For the rest of this section specialize \(n=a\in\{5,7\}\), in
characteristic different from \(3,a\). Thus \(N\) is the proper
geometric pro-three group of \(C_a\). These are the
base exponents whose affine boundary problem remains open.


The diagonal group \[
 A_a=(\mu_a)^2\subset\operatorname{Aut}(C_{a,\bar k})
\] acts faithfully on \(H_1(C_{a,\bar k},\mathbb F_3)=N/\Phi(N)\).

Indeed, every nonidentity element \(\gamma\in A_a\) has order \(a\). The quotient map \(C_a\to C_a/\langle\gamma\rangle\) is tame, since the residue characteristic is not \(a\). Riemann--Hurwitz gives \[
 g(C_a/\langle\gamma\rangle)<g(C_a).
\] The invariant subspace of rational three-adic \(H_1\) has dimension twice the quotient genus, so \(\gamma\) acts nontrivially on \(H_1(C_a,\mathbb Q_3)\).

If its action modulo three were trivial, its integral matrix would lie in \(\ker(\operatorname{GL}_{2g}(\mathbb Z_3)\to
       \operatorname{GL}_{2g}(\mathbb F_3))\), a pro-three group. A nontrivial element of order \(a\) cannot lie there. Therefore its residual action is nontrivial. Every nonidentity element of \(A_a\) has order \(a\), proving faithfulness.

The criterion of Lemma~\ref{lem:ff3-trivial-residual} therefore does not construct the same-component boundary collision sought here. Failure of a sufficient criterion is not proof that such a collision cannot occur.

\subsection{Hoshi’s general counterexample mechanism and the boundary field}
\label{subsec:ff3-hoshi}

\citet[Remark~10(i), pp.~742--743]{FCHoshi2009} explains a general source of collisions in proper geometrically pro-\(\ell\) point maps over finite fields. One first makes the outer Galois action pro-\(\ell\) by a finite extension of constants; one then uses a further extension of degree prime to \(\ell\) that supplies new points. Conjugate points in that extension yield coincident section classes in the stated setting.

Hoshi writes conjugacy by the arithmetic group. For sections of the full procyclic \(G_k\), this is equivalent to geometric conjugacy: a conjugator with arithmetic image \(g\) can be multiplied on the right by \(s(g)^{-1}\), which centralizes \(s(G_k)\), to obtain a conjugator in the geometric kernel.

This is a genuine warning against unrestricted singleton-prime point injectivity. It does not specify the Fermat boundary or the fibre above a rational coordinate point.

For the present boundary, the mechanism has a concrete obstruction.

\begin{proposition}[The boundary field in the constant-extension construction]
\label{prop:ff3-hoshi-boundary}

Let \(k'=\mathbb F_{Q^t}\) be a finite constant-field extension over which the outer \(G_{k'}\)-action on \(N\) has pro-three image. Then \[
 \operatorname{ord}_a(Q^t)
\] is a power of three. Any further extension of \(k'\) of degree prime to three adds no points to the coordinate boundary \(D(k')\).

\end{proposition}
\begin{proof}

The residual action of \(G_{k'}\) on \(N/\Phi(N)\) is a finite three-group. Thus its Frobenius matrix has three-power order.

The faithful diagonal subgroup \(A_a\) from Section~\ref{subsec:ff3-automorphisms} is normalized by Frobenius. On \(A_a\), conjugation is the \(Q^t\)-power action on its root-of-unity coordinates (up to replacing Frobenius by its inverse, which does not change the order). Faithfulness on the residual module therefore identifies this conjugation with conjugation by the residual Frobenius matrix. Its order must be a power of three. The order of scalar \(Q^t\) on \((\mathbb Z/a\mathbb Z)^2\) is exactly \(\operatorname{ord}_a(Q^t)\), proving the first assertion.

The field of all coordinate points over \(k'\) is \(k'(\mu_a)\), of degree \(\operatorname{ord}_a(Q^t)\). For \(a=5\), this degree divides four, hence must already be one. For \(a=7\), it divides six, hence is one or three. More generally its three-power part cannot be removed by an extension of degree prime to three: \[
 \operatorname{ord}_a(Q^{tu})
 =\frac{\operatorname{ord}_a(Q^t)}
        {\gcd(\operatorname{ord}_a(Q^t),u)}
 =\operatorname{ord}_a(Q^t)
 \qquad(3\nmid u).
\] Since \(a\) is prime, a coordinate component has either its single canonical rational point or all \(a\) points. Thus the further extension supplies no new coordinate boundary points.

\end{proof}

This excludes this particular use of Hoshi's two-stage construction for producing new Fermat boundary points. It does not decide the same-component alternative over a field where all boundary points are already rational, nor does it exclude other collision constructions.

\subsection{A depth-three limitation at quintic primes}
\label{subsec:ff3-quintic-depth-three}

The full pro-three boundary theorem cannot in general be detected
in the first three ordinary lower-central grades.  The following
result gives a specified positive-density family and an explicit prime.

\begin{proposition}[Constant quintic point map through class three]
\label{prop:ff3-quintic-depth-three}
Let $q\ne3,5$ be a rational prime with $q\equiv1\pmod5$, let
\[
 C/\mathbb F_q:\ X^5+Y^5+Z^5=0,\qquad
 C_A:\ v^2+v=u^5,\qquad A=\operatorname{Jac}(C_A),
\]
and suppose arithmetic Frobenius on $A[3]$ has order twenty.
Use $b=(0:-1:1)$ as base point and put
\[
 N=\pi_1(C_{\overline{\mathbb F}_q},b)^{\mathrm{pro}\text{-}3}.
\]
With closed ordinary lower-central terms, continuous nonabelian
cohomology satisfies
\[
 H^1(G_{\mathbb F_q},N/\gamma_4N)=\{1\}.
\]
Thus every section of this arithmetic quotient is geometrically
conjugate to the point section of $b$.  In particular, the proper
point-section map is constant in this quotient and in every finite
Galois-stable quotient of $N$ of nilpotency class at most three.

The set of rational primes satisfying these conditions has natural
density $1/10$, and it contains $q=11$.
\end{proposition}

\begin{proof}
The exact residual calculation for the quintic CM Jacobian gives
\[
 E=\mathbb Q(\zeta_5)\subset K=\mathbb Q(A[3]),\qquad
 [K:\mathbb Q]=80,\qquad
 \operatorname{Gal}(K/E)=C_{20}\subset\mathbb F_{81}^{\times};
\]
see Lemma~\ref{lem:lm-quintic-residual-group}.
Here $A[3]$ has rank one over
$\mathbb Z[\zeta_5]/3=\mathbb F_{81}$.
The three degree-five quotient maps
\[
 \begin{split}
 f_Z&=(-XY/Z^2,X^5/Z^5),\\
 f_X&=(-YZ/X^2,Y^5/X^5),\\
 f_Y&=(-ZX/Y^2,Z^5/Y^5)
 \end{split}
\]
give the integral isomorphism
\[
 T_3\operatorname{Jac}(X_5)\simeq(T_3A)^3.
\]
Indeed their joint pushforward $\Phi$ and summed pullback $\Psi$
satisfy $\Phi\Psi=[5]$ and $\Psi\Phi=[5]$, by the three disjoint
diagonal-character summands.  Hence $5^{-1}\Psi$ is an integral
three-adic inverse.  Good reduction at $q$ transports this identity
to the finite-field Frobenius representation.

Because $q$ splits in $E$, its residual Frobenius is multiplication
by an element $t\in\mathbb F_{81}^{\times}$ of order twenty.
The eigenvalues over $\overline{\mathbb F}_3$ on $A[3]$ are
\[
 t,\quad t^3,\quad t^9,\quad t^{27}=t^7.
\]
Each occurs three times on
$\overline V=N^{\mathrm{ab}}/3$.
The action is semisimple since twenty is prime to three.
None of the exponents $1,3,7,9$ is zero modulo twenty;
every sum of two lies between two and eighteen;
every sum of three is odd.
Consequently $1$ is not an eigenvalue of Frobenius on
\[
 \overline V^{\otimes m},\qquad 1\le m\le3.
\]
Repeated tensor factors are included.  Replacing $t$ by another
generator, or Frobenius by its inverse, does not change this assertion.

Put $L_m=\gamma_mN/\gamma_{m+1}N$.
Iterated commutators give a continuous equivariant surjection
\[
 (N^{\mathrm{ab}})^{\otimes m}\longrightarrow L_m;
\]
the grades are finitely generated compact $\mathbb Z_3$-modules.
The commutator description and its compatibility with the Galois
action are also described in
\citet[Introduction, author pp.~1--2]{DavisPriesWickelgren2024}.
It follows that $F-1$ is invertible modulo three on $L_m$ for
$m\le3$. Nakayama's lemma makes it surjective integrally, and a
surjective endomorphism of a finitely generated $\mathbb Z_3$-module
is injective. Thus
\[
 H^0(G_{\mathbb F_q},L_m)=H^1(G_{\mathbb F_q},L_m)=0
 \qquad(1\le m\le3).
\]
For $H^1$, explicitly choose a coboundary with the prescribed
cocycle value at the topological Frobenius generator; equality on
the full procyclic group follows by continuity.  The arithmetic
quotient remains full $G_{\mathbb F_q}$.

Induct through the central extensions
\[
 1\longrightarrow L_m\longrightarrow N/\gamma_{m+1}N
 \longrightarrow N/\gamma_mN\longrightarrow1.
\]
The abelian case follows from the displayed vanishing.  In a
successive stage, conjugate a cocycle so that its image in the
preceding quotient is trivial, lifting one conjugating element
through the surjection.  The remaining cocycle takes values in
$L_m$ and is a coboundary.  This proves the asserted singleton
cohomology set through $m=3$.  Equivalently, an invariant path lifts
at each central step because its fibre is an $L_m$-torsor with
zero $H^1$; the zero $H^0$ also makes these low-depth paths unique.
No section in the full geometric group is assumed to exist above
an arbitrary quotient section.

The residual elements defining the prime set are the eight generators
of the normal subgroup $C_{20}$ in the group of order eighty.
This is a conjugacy-stable subset, so Chebotarev gives natural
density $8/80=1/10$.

For an explicit prime, exact counting on $v^2+v=u^5$ gives
\[
 \#C_A(\mathbb F_{11})=13,\qquad
 \#C_A(\mathbb F_{121})=103.
\]
One can use $\mathbb F_{121}=\mathbb F_{11}[s]/(s^2-2)$,
where $2$ is a nonsquare, and count the square roots of
$1+4u^5$.  The Frobenius polynomial is therefore
\[
 T^4+T^3-9T^2+11T+121.
\]
Modulo three it is $T^4+T^3+2T+1$.
Its derivative is $(T-1)^3$ and its value at one is two, so it is
square-free.  Its companion matrix has order twenty by exact
matrix multiplication over $\mathbb F_3$.  Thus $q=11$ meets
the hypothesis.  These finite calculations are recorded in the
reproducibility certificate.
\end{proof}

At $q=11$, the same isogeny gives $\#C(\mathbb F_{11})=15$,
so all proper rational points are the fifteen coordinate points.
Their equality in $N/\gamma_4N$ is compatible with
Theorem~\ref{thm:ff3-different-components}: the latter concerns
the full geometric pro-three group and separates different
coordinate components. The proposition makes no assertion of a
full collision within one component. Already in tensor degree four,
$1+1+9+9\equiv0\pmod{20}$, so the eigenvalue exclusion stops.
An actual integral higher-path class remains to be determined;
this limitation of low-depth quotients does not settle the affine
base-five or base-seven Selmer-section problem.

\subsection{Two obstructions to diagonal-symmetry shortcuts}
\label{subsec:ff3-diagonal-shortcuts}

The next two observations explain why coprime complement conjugacy
does not settle the remaining full same-component boundary question.
They concern the entire Fermat geometric group, not just one
hyperelliptic quotient.

\begin{proposition}[A nontrivial diagonal twist cannot remove the residual obstruction]
\label{prop:ff3-diagonal-twist-obstruction}
Let $a\in\{5,7\}$, let $k=\mathbb F_Q$ have characteristic prime
to $3a$, and assume $\mu_a\subset k$.  On
$C_a:X^a+Y^a+Z^a=0$, let
\[
 \delta(X:Y:Z)=(\zeta^rX:\zeta^sY:Z),\qquad
 r,s\in\mathbb Z/a\mathbb Z,
\]
where $\zeta$ is primitive.  Let $N$ be the geometric pro-three
fundamental group and let $F$ denote arithmetic Frobenius.
If the action of $\delta^{-1}F$ on $N/\Phi(N)$ has three-power
order, then $r=s=0$.
In particular, for a nonidentity $\delta$, the closed subgroup of
$\operatorname{Out}(N)$ generated by its twisted Frobenius cannot
be pro-three.
\end{proposition}

\begin{proof}
Here $\delta$ and $F$ are initially outer automorphisms of $N$;
their actions on $N/\Phi(N)$ are canonical. They commute as outer
automorphisms because $\delta$ is defined over $k$. Thus the
displayed twisted action is defined without choosing a
$\delta$-fixed base point.

Put $J_a=\operatorname{Jac}(v^2+v=u^a)$, and write $\beta$ for
the automorphism induced by $(u,v)\mapsto(\zeta u,v)$.
The three degree-$a$ morphisms to this same curve are
\[
 \begin{split}
 f_Z&=(-XY/Z^2,X^a/Z^a),\\
 f_X&=(-YZ/X^2,Y^a/X^a),\\
 f_Y&=(-ZX/Y^2,Z^a/Y^a).
 \end{split}
\]
They satisfy $u^a=v(v+1)$ and extend to the proper smooth models.
Each is the quotient by its order-$a$ diagonal deck subgroup.
The induced pushforward on three-torsion is surjective: its
composition with pullback is multiplication by $a$, which is
invertible on three-torsion.

The three $u$-coordinate weights of $\delta$ are
\[
 w_Z=r+s,\qquad w_X=s-2r,\qquad w_Y=r-2s
 \pmod a.
\]
Consequently these quotient maps intertwine the action of
$\delta$ with $\beta^{w_i}$ and the action of $F$ with the same
Frobenius operator $B$ on $J_a[3]$.

The quotient of $v^2+v=u^a$ by $\beta$ is the projective $v$-line.
The $\beta$-invariants in its rational first homology therefore
vanish. Since $a$ is invertible in $\mathbb Z_3$, invariants are
an integral direct summand and commute with reduction modulo three.
Moreover $\dim_{\mathbb F_3}J_a[3]=a-1$, and $\Phi_a(T)$ is
irreducible over $\mathbb F_3$: the order of $3$ modulo $a$ is
$a-1$ for both $a=5$ and $a=7$. Thus $J_a[3]$ has rank one over
\[
 \mathbb F_3[\beta]\simeq\mathbb F_{3^{a-1}}.
\]
Frobenius commutes with $\beta$ because $\zeta\in k$, so $B$ is
a nonzero scalar in this field.

The three quotient actions of $\delta^{-1}F$ are therefore the
scalars $B\beta^{-w_i}$. Under the hypothesis each has
three-power order. Its order also divides $3^{a-1}-1$, which is
prime to three, so each is the identity. Hence
\[
 \beta^{w_Z}=\beta^{w_X}=\beta^{w_Y}.
\]
As $\beta$ has exact order $a$, this gives $3r=3s=0$ modulo $a$,
and therefore $r=s=0$. The final assertion follows because the
finite residual image of a pro-three group is a three-group.
\end{proof}

This excludes a nontrivial diagonal twist as the first step in the
pro-three-outer-action constant-extension construction. It does not
exclude other constructions of collisions, or decide the fibres
over a field where all boundary points are already rational.

\begin{proposition}[The geometric diagonal centralizer has a nonzero quotient]
\label{prop:ff3-diagonal-centralizer}
Let $a\in\{5,7\}$ and let $C_a$ be the proper Fermat curve over
an algebraically closed field of characteristic prime to $3a$.
Let $N$ be its geometric pro-three fundamental group and
$H=(\mu_a)^2$ its diagonal automorphism group. Form the relative
geometric action extension
\[
 1\longrightarrow N\longrightarrow E_H\longrightarrow H
 \longrightarrow1,
\]
retaining the full finite quotient $H$. Choose any complement
$S\subset E_H$, mapping isomorphically onto $H$, and let it act
on $N$ by conjugation. Such a complement exists by coprime
Schur--Zassenhaus. Put $L_2=\gamma_2N/\gamma_3N$ and
$g=(a-1)(a-2)/2$. Then
\[
 Z_N(S)\subset\gamma_2N,\qquad
 \frac{Z_N(S)}{Z_N(S)\cap\gamma_3N}
       \simeq L_2^H\simeq\mathbb Z_3^{g-1}.
\]
In particular this quotient has rank five for $a=5$ and rank
fourteen for $a=7$.
\end{proposition}

\begin{proof}
The choice of complement supplies an actual action on $N$.
Its action on each ordinary lower-central grade is the canonical
geometric $H$-action: changing lifts by inner automorphisms has
no effect on a grade. No arithmetic Galois invariance of $S$ is
assumed or asserted.

Write $V=N^{\mathrm{ab}}$. After an unramified coefficient
extension of $\mathbb Z_3$ containing $\mu_a$, the characters of
$H$ on $V$ are
\[
 (i,j)\in(\mathbb Z/a\mathbb Z)^2,
 \qquad i\ne0,\quad j\ne0,\quad i+j\ne0,
\]
each with multiplicity one. In characteristic zero, the
differentials $x^{i-1}y^{j-a}\,dx$, with $i,j\ge1$ and
$i+j\le a-1$, have characters $(i,j)$; adjoining their conjugates
gives precisely the displayed set. Smooth proper
prime-to-characteristic specialization gives the same geometric
representation in the stated characteristics. The degree-two
surface description used below is also given in
\citet[author pp.~2--4]{DavisPriesWickelgren2024}.
In particular $V^H=0$, so every element of $Z_N(S)$ lies in
$\gamma_2N$.

The degree-two surface presentation gives
\[
 L_2\simeq\bigwedge^2V/\mathbb Z_3\omega,
\]
where $\omega$ is the primitive symplectic relation. A trivial
diagonal character in the exterior square pairs $(i,j)$ with
$(-i,-j)$. There are exactly $g$ such pairs. The relation is
invariant and primitive, with a unit coefficient in a symplectic
basis. Since $|H|=a^2$ is invertible in $\mathbb Z_3$, the
averaging projector makes invariants exact. Therefore taking
invariants removes exactly one primitive rank-one summand,
and $L_2^H$ is finite free of rank $g-1$. The unramified
coefficient extension is faithfully flat, so this rank and
freeness assertion descends to $\mathbb Z_3$.

It remains to check that these graded fixed elements actually
lift to the geometric centralizer. For
$\overline x\in L_2^H$, choose $x\in\gamma_2N$.
The elements
\[
 c_h=x^{-1}h(x)\in\gamma_3N\qquad(h\in H)
\]
form a continuous nonabelian cocycle for the action furnished by
$S$. Its class is trivial because $H$ is finite of order prime
to three and $\gamma_3N$ is pro-three. Explicitly, finite
Schur--Zassenhaus conjugacy kills it in every finite
$H$-stable quotient; the corresponding nonempty closed sets of
conjugating elements have the finite-intersection property, so
compactness kills it in $\gamma_3N$ itself. Hence there is
$u\in\gamma_3N$ with $c_h=u\,h(u)^{-1}$ for every $h$.
Then $xu$ is fixed by $H$ and lifts $\overline x$.
Thus $Z_N(S)\to L_2^H$ is surjective, with kernel
$Z_N(S)\cap\gamma_3N$, proving the stated quotient.
\end{proof}

The absence of diagonal invariants in the original Jacobian does
not make this geometric centralizer trivial. Coprime conjugacy
of complements therefore does not provide the Frobenius
compatibility required for a rational invariant path. Neither
proposition proves a full same-component collision or full
separation, and neither supplies the missing affine
base-five or base-seven Selmer-section classification.


\subsection{Exact remaining gate}
\label{subsec:ff3-gate}

For \(Q\equiv1\pmod a\), choose a primitive \(\zeta\in k\), put \[
 b=(0:1:-1),\qquad
 \delta(X:Y:Z)=(X:Y:\zeta Z).
\] The remaining question is exactly whether \begin{equation}
\mathcal P_3(b,\delta b)^{G_k}
\label{eq:ff3-11-1}
\end{equation} is empty. By \eqref{eq:ff3-3-1}, nonemptiness is equivalent to \(\kappa(b)=\kappa(\delta b)\). Since \(a\) is prime, that one equality implies the collapse of the whole coordinate component, and then of each component by permutation symmetry. Theorem~\ref{thm:ff3-different-components} still separates the three resulting classes.

One can express the same gate as a nonabelian Lang equation. Choose a path \(p\) and use the right \(N\)-torsor convention \(g(p)=p\,c_g\). Then \(c_g\) is a continuous cocycle. For arithmetic Frobenius \(F\), a path \(pn\) is fixed if and only if \begin{equation}
c_F F(n)=n,\qquad\text{equivalently}\qquad
 c_F=nF(n)^{-1}.
\label{eq:ff3-11-2}
\end{equation} The value at a topological generator determines invariance under \(G_k\), by continuity. The abelianization of this equation is always soluble for the boundary difference, as Section~\ref{subsec:ff3-abelian} explains. Full integral pro-three solvability is the unresolved issue.

No specific solution or obstruction to \eqref{eq:ff3-11-1} at split rational primes has been obtained in this proof. Consequently the following claims remain unproved here:

\begin{itemize}
\item
  that another boundary point lies over a chosen rational cusp for even one specified split rational prime;
\item
  that collisions occur at a set of positive Dirichlet density;
\item
  that singleton boundary fibres occur at a density-one set;
\item
  that the affine pro-three base-five or base-seven Selmer section classification follows from this partial result.
\end{itemize}

The cofinite-prime almost-injectivity theorems do not fill this gap. Saïdi--Tamagawa Proposition 1.8(vi)--(vii) and Corollary 1.11(iii), printed pp.~144--148, impose a cofinite retained prime set. The set \(\{3\}\) is not cofinite. Their earlier Lemma 1.7, used in Section~\ref{subsec:ff3-centralizers}, is valid in the singleton setting and supplies only the centralizer input stated there.
